\chapter*{Abstract}

The project \textbf{"Better Spotify / Sample Browser"} (hereinafter referred to as \textbf{Sample Browser}) was developed by students Bogdan Glavan, Andrei Simon, Nikita Creciun, Leonid Tiscenco, and Mihai-Constantin S\^{i}rghii, under the mentorship of Prof. Elena Graur at the Technical University of Moldova, Faculty of Computers, Informatics and Microelectronics. This report encapsulates the initial Conceptual Design of a Software Application (PCAS) stage, covering problem definition, domain analysis, methodological evolution, solution architecture, target audience validation, and comparative benchmarking.

\textbf{Keywords: } Better Spotify, Sample Browser, Audio Sampling, Music Information Retrieval (MIR), Granular Passage Discovery, Music Streaming Architectures, Creator Workflow, Problem-Based Learning.

The primary objective of \textbf{Sample Browser} is to eliminate the persistent workflow friction experienced by independent music producers during sample crate-digging across streaming catalogs. Current digital streaming platforms index musical compositions exclusively at the whole-track metadata level, ignoring sub-track time boundaries. Consequently, producers seeking specific localized musical textures---such as isolated 4-bar rhythm breaks or distinct instrumental passages---must scrub back and forth manually across long audio files, disrupting their creative focus. To solve this, \textbf{Sample Browser} provides a specialized discovery platform that interfaces with embedded streaming audio via the YouTube IFrame API to segment, tag, and recommend localized audio passages (4--30 seconds). By coupling a structured Two-Tier Tagging System (objective structural presets and contextual descriptors) with direct timestamp seeking and automatic loop auditioning, the system transforms unstructured listening into a targeted creator utility. Through a rigorous Problem-Based Learning methodology, this report documents the team's pivot from an unfeasible consumer streaming overlay, establishes four verifiable validation hypotheses (H1--H4), benchmarks market alternatives, and provides the architectural foundation for subsequent systems modeling and empirical user validation.
